problem:  
Let \( (M^{2n}, J, g) \) be a compact complex manifold endowed with a strong Kähler with torsion (SKT) metric, meaning its fundamental form \( \omega \) satisfies \( \partial\bar{\partial}\omega = 0 \). Consider the **pluriclosed flow**
\[
\frac{\partial \omega_t}{\partial t} = -(\rho^{B})^{1,1},
\]
where \( \rho^{B} \) is the Bismut–Ricci form of the Hermitian metric \( g_t \).  

Assume that:
1. The first Chern class \( c_1(M) = 0 \).
2. The Bott–Chern cohomology class \( [\omega_0]_{BC} \) lies in the closure of the cone of classes representable by positive \( (1,1) \)-currents (so that the flow begins in a topologically "positive" direction).  

**Question:**  
Does the pluriclosed flow starting from such an SKT metric \( \omega_0 \) exist for all time and converge, up to diffeomorphism, to a pluriclosed Hermitian metric \( \omega_\infty \) whose Bismut–Ricci form vanishes?  
If so, does this limiting condition imply that the Bismut connection of \( g_\infty \) has holonomy contained in \( SU(n) \), thereby corresponding to a "Calabi–Yau structure with torsion" in the SKT category?  

Why is it a "good" problem:  
This problem aims to identify whether pluriclosed flow—an analogue of the Kähler–Ricci flow in the non‑Kähler setting—can dynamically produce torsion Calabi–Yau structures under minimal geometric assumptions. It generalizes the Calabi–Yau theorem and the long‑time behavior of Kähler–Ricci flow to the SKT realm, where torsion plays a fundamental role.  

The question is naturally motivated by existing partial results (notably by Streets–Tian, Ustinovskiy, and others), but remains open in general dimensions and without symmetry or surface restrictions. The assumption \(c_1(M)=0\) aligns with the topological Calabi–Yau condition, and the Bott–Chern positivity captures a plausible non‑Kähler analog of the Kähler class positivity used in Ricci flow theory.  

If resolved, a positive answer would establish a powerful non‑Kähler version of the Calabi–Yau paradigm, linking analytic convergence under pluriclosed flow to Bismut holonomy reductions—bridging geometric analysis, complex Hermitian geometry, and the study of torsion structures in generalized complex geometry.  

Its statement is compact, geometric, and natural, yet the problem probes deeply into uncharted territory of complex differential geometry and geometric PDEs, satisfying the essential qualities of a “good” research problem: profound, cross‑disciplinary, and conceptually simple but mathematically formidable.